\subsection{Delineation choice: NHDPlus HR versus threshold-based DEM delineation}
\label{sec:methods-delineation}

The hydrography preprocessing of \S\ref{sec:methods-nhd} conditions delineation on the NHDPlus HR existing network. QSWAT\textsuperscript{+} \citep{QSWATplusSoftware} also supports a fully DEM-based alternative in which the channel network is extracted from flow accumulation above a support-area threshold using TauDEM \citep{Tarboton1997TauDEM}, a long-recognized sensitivity in channel extraction whose subdivision level propagates into SWAT-simulated flow and loads \citep{Tarboton1991ChannelNetworks, Jha2004Subdivision, Chaplot2005DEMmesh}. We therefore ask whether conditioning on NHDPlus HR is necessary or whether standard threshold delineation would yield an equivalent model, and address it with a controlled, same-DEM comparison on the small (S) benchmark basin and a watershed-scale build of the Peace River HUC-8 (full analysis, figures, and the Silver-River karst case in Supplementary Section~\ref{sec:supp-delineation}).

Two properties favor NHDPlus-HR conditioning, neither a claim that threshold delineation is incapable. First, \emph{determinism}: the conditioned contributing area is fixed by the surveyed, hydro-enforced hydrography \citep{Moore2019NHDPlusHR, Lindsay2016StreamBurning}, whereas the threshold area is extent-dependent---on the S basin it brackets but does not reproduce the conditioned value (square box $67.6$, basin-clipped $43.1$, versus $52.6$~km\textsuperscript{2}). Second, \emph{correctness without per-basin tuning}: NHDPlus HR wires the basin's lakes as reservoir objects directly, whereas the DEM alternative required a specific lake method and threshold on the S basin and, on the lake-dense Peace HUC-8 ($347$ waterbodies), could be made to build only by clipping the DEM, selectively burning the major rivers, and dropping the lakes---yielding a heavily over-segmented network ($2{,}304$ subbasins vs.\ $162$). For an unattended generator across thousands of domains these are the operative advantages; the comparison is explorable on the companion page \citep{SWATGenXcompanion}.
